ปัญหาภาวะสารสนเทศท่วมท้นและการรายงานข่าวซ้ำซ้อนจากหลากหลายแหล่ง ทำให้ผู้อ่านต้องเสียเวลาในการคัดกรองและติดตามข้อมูล โครงงานนี้จึงมีวัตถุประสงค์เพื่อศึกษา ออกแบบ และพัฒนาระบบรวบรวม จำแนกหมวดหมู่ และวิเคราะห์ข่าวสารอัตโนมัติแบบครบวงจร (End-to-End News Pipeline) โดยสถาปัตยกรรมระบบได้รับการออกแบบตามหลักการ SOLID และ GRASP เชื่อมต่อการทำงานเป็นเลเยอร์อย่างเป็นระเบียบ ส่วนบริการเบื้องหลังพัฒนาด้วยภาษา Python และ FastAPI ร่วมกับ Playwright Headless Browser ในการดึงข้อมูลข่าวสารจาก 14 สำนักข่าวชั้นนำของไทย พร้อมระบบแคชหลายระดับเพื่อลดภาระการเชื่อมต่อเครือข่าย ระบบนำเสนอไปป์ไลน์การประมวลผลภาษาธรรมชาติที่ครอบคลุม การตรวจจับและจัดกลุ่มข่าวที่เป็นเหตุการณ์เดียวกัน (Story Clustering) เพื่อลดความซ้ำซ้อน, การจัดอันดับกระแสความนิยมด้วยอัลกอริทึมเชิงสถิติ (Deterministic Trending Score) ที่ลดทอนข่าวประจำวันโดยไม่พึ่งพาโมเดลภาษาขนาดใหญ่, การสังเคราะห์สรุปเนื้อหาและวิเคราะห์อารมณ์ข่าวด้วย KKU Gen.AI API โดยแสดงผลพรีวิวบางส่วนควบคู่กับลิงก์อ้างอิงต้นฉบับเพื่อคุ้มครองลิขสิทธิ์ และการจำแนกเนื้อหา 9 หมวดหมู่ด้วยแนวทางไฮบริดผสานกฎคำสำคัญเชิงบริบทเข้ากับโมเดลการเรียนรู้ของเครื่อง สำหรับส่วนติดต่อผู้ใช้พัฒนาด้วย Vanilla JavaScript ในรูปแบบ Single Page Application (SPA) รับส่งข้อมูลแบบเรียลไทม์ผ่าน WebSocket (Socket.IO) พร้อมระบบแนะนำข่าวสารเฉพาะบุคคลที่ประมวลผลบนฝั่งผู้ใช้งานเพื่อรักษาความเป็นส่วนตัวตามมาตรฐาน PDPA (Zero-PII) จากการทดสอบระบบผ่านชุดทดสอบอัตโนมัติทั้ง 6 มิติ รวม 423 ชุดทดสอบ (ร้อยละ 100) และผลการประเมินความพึงพอใจจากกลุ่มตัวอย่างผู้ใช้งานจริงพบว่าอยู่ในระดับดีมาก โดยมีคะแนนเฉลี่ยรวม 4.86 จาก 5.00 คะแนน และมีแนวโน้มการแนะนำบอกต่อร้อยละ 90.43 ชี้ให้เห็นว่าระบบช่วยประหยัดเวลาและเพิ่มประสิทธิภาพในการติดตามข่าวสารได้อย่างเป็นรูปธรรม
